\documentclass[fontsize=12pt]{mybib}

\setheaderlogo{mnr-green}

\newcommand{\place}{ \\\\\\\\\\\\\\\\\\\\\\ }

% ensure scrlayer-scrpage has sufficient footheight
\renewcommand{\DFA}{12mm}
\renewcommand{\DZA}{1.2}
\setlength{\footheight}{30.4pt}
\author{OTS}
\begin{document}

\section{Johannes 1}

\subsection{Einleitung}
Im Johannes-Evangelium wird der Autor nirgends genannt. Irenäus (130 - 200 \nChr) erwähnt, dass Polykarb (ca. 69 - 155 \nChr) sagte, dass Johannes im fortgeschrittenen Alter das Evangelium verfasst habe.\\
Clemens von Alexandria (ca. 150 - 215 \nChr) erwähnte, dass Johannes im Bewusstsein der anderen Evangelien und im Wirken des hl. Geistes sein Evangelium abgefasst hatte.\\
Während in den anderen Evangelien Johnnes der Apostel mehr als 20mal erwähnt wird, wird der Name Johannes vom Apostel nie erwähnt. Johannes erwähnt sich nur indirekt in \biblezit{Joh}{13:23}, \biblezit{Joh}{19:26} und \biblezit{Joh}{21:2};. \\
Geschrieben wurde das Evangelium vermutlich zwischen 90 und 100 n.Chr. in Ephesus.

\subsection{ Meine Fragen }
\begin{enumerate}
    \item Was meint Johannes mit: \enquote{Im Anfang war das Wort}?\place
    \item Von welchen Johannes wird im Vers 6 gesprochen? Wo erfahren wir mehr über diesen Johannes? Was sagt dieser Johannes über sich, wer er ist?\place
    \item Ab Vers 9 - 13 wird über das Licht gesprochen. Wieso nicht mehr vom Wort?\place
\end{enumerate}

\subsection{Fragen von Wiersby}
\begin{enumerate}
    \item Wie wird Jesus heute von der Welt gesehen?\place
     \item Wie hat Johannes in diesem ersten Kapitel Jesus geschildert?\place
    \item Was sagt uns jeder von Johannes verwendeter Name und Titel über Jesus?\place
    \item Warum ist das Wort Fleisch geworden?\place
    \item Was erfahren Sie in diesem ersten Kapitel über Johannes den Täufer und seine Beziehung zu Jesus?\place
    \item Als der sechste Tag seines öffentlichen Wirkens zu Ende ging, hatte Jesus sechs Jünger. Wie haben diese Männer zum Glauben gefunden?\place
    \item Was können Sie über Evangelistion lernen, wenn Sie die Art uns Weise betrachten, wie diese Männer zuz Jüngeren Jesu wurden?\place
    \item Was ist für Sie das Wichtigste, das Sie in diesem Kapitel über Jesus gelernt haben?\place
    \item Welche Auswirkungen hat das in dieser Woche Gelernte auf Ihren Alltag?\place
\end{enumerate}

\end{document}